\section{Related Work}
\label{sec:related-work}
%\todo[inline]{We acknowledge that techniques for
%sequential rule mining are related.
%Yet, there is a fundamental difference in the adopted
%query model.
%In the revised version, we will add a more
%comprehensive discussion of these techniques,
%clarifying their differences and the unique
%challenges in the context of CEP.}
Our work relates to existing approaches for automated query discovery, as
well as incremental algorithms for sequential pattern mining and frequent
itemset mining.

\textbf{Automated Query Discovery}
So far research in the area of automated query discovery did not take into account
updates in the database or the support threshold.
We adopted the query model used in ~\cite{DBLP:conf/btw/Kleest-Meissner23}. But
the proposed algorithm does not find all matching queries but only discovers
one matching query for each run.
The other approaches \cite{ilminer,icep} further lack in the query
expressiveness. While \cite{icep} cannot discover queries with
multiple occurrences of an attribute value, \cite{ilminer} lacks
the ability to find query terms which only contain variables.

\textbf{Incremental Sequential Pattern Mining}
A related field in which incremental algorithms have been proposed
is sequential pattern mining. Here, the input is a database of strings
rather than a stream of events which can contain multiple attributes.
Furthermore, the discovery is limited to type level without the
ability to discover variables or patterns.
In \cite{zheng2002,zheng2003}, the authors propose two algorithms, one for
adding
and one for deleting data to the original database. The algorithms make use of
negative border sequences which are those sequences whose sub-sequences
are frequent to recalculate the results.
\cite{mallick2013} presents a survey that summarizes the incremental
approaches in this
area. There the highlighted algorithms are based on common frequent
sequence mining algorithms.





\textbf{Incremental Frequent Itemset Mining}
In frequent itemset mining sets of items are discovered that occur frequently
together in a database. This setting differs from our query discovery since
the events in a discovered query are order sensitive whereas item sets can occur
in arbitrary order.
High-utility pattern mining, a sub-field of frequent itemset mining,
discovers itemsets based on a user-defined utility threshold.
The incremental algorithm introduced in \cite{yun2015} proposes a tree
structure to update the item sets.

\textbf{Sequential Rule Mining}
Sequential rule mining \cite{fournier2011rulegrowth,erminer} is a combination of sequential pattern mining and
frequent itemset mining. The goal is to find rules that describe a
relationship between itemsets in a sequence. The rules are of the form
$X \rightarrow Y$ where $X$ and $Y$ are itemsets. The rules are
evaluated based on metrics like support and
confidence, and my also be found
incrementally~\cite{ierminer}.
However, one limitation of such rules is that they are
bounded to
one implication whereas our queries can contain multiple events.
Also, the rules are limited in their expressiveness as
the itemset captures only type-level information, not
allowing for variables.



% Theory paper recommended by Sarah
% \cite{idris2020}